% !TEX root = ../main.tex
\section{Решаемые проблемы и практическая значимость}
\addcontentsline{toc}{section}{Решаемые проблемы и практическая значимость}

\subsection{Проблема запоминания и применения диапазонов рук}
Одной из фундаментальных сложностей в изучении покерной стратегии является необходимость запоминания и корректного применения префлоп-диапазонов для различных игровых ситуаций. Диапазон (range) --- это набор стартовых рук, с которыми игрок должен предпринимать определённое действие (рейз, колл, фолд) в зависимости от:

\begin{itemize}
    \item своей позиции за столом (баттон, кат-офф, биг-блайнд и др.),
    \item предыдущих действий оппонентов (открытие рейзом, 3-бет и т.д.),
    \item типа игры (кэш-игра, турнир с глубокими или короткими стеками).
\end{itemize}

Современные покерные стратегии предполагают использование десятков различных диапазонов, каждый из которых может содержать до 169 уникальных комбинаций стартовых рук (с учётом мастей и рангов). Запоминание и корректное применение таких объёмов информации в условиях ограниченного времени на принятие решения (особенно в онлайн-покере) представляет значительную когнитивную нагрузку даже для опытных игроков.

Традиционный подход --- использование распечатанных чартов диапазонов, размещаемых рядом с монитором. Однако данный метод имеет существенные недостатки:

\begin{itemize}
    \item Необходимость постоянного переключения внимания между игровым окном и физическим чартом, что снижает концентрацию и скорость реакции.
    
    \item Отсутствие динамической адаптации: чарт не учитывает текущую игровую ситуацию (размеры стеков, позиции активных игроков), требуя от игрока дополнительных ментальных вычислений.
    
    \item Невозможность быстрой смены стратегии: при переходе между различными типами игр (кэш $\rightarrow$ MTT) требуется замена физических материалов.
\end{itemize}

\subsection{Практическая значимость разработанного решения}
Разрабатываемое программное обеспечение решает указанные проблемы путём интеграции стратегических диапазонов непосредственно в процесс игры:

\begin{itemize}
    \item \textbf{Автоматизация контроля соблюдения диапазонов:} система в реальном времени анализирует текущую ситуацию, сопоставляет её с загруженными пользователем диапазонами и визуально подсвечивает корректное действие, исключая ошибки, связанные с выходом за рамки стратегии.
    
    \item \textbf{Снижение когнитивной нагрузки:} игроку не требуется запоминать десятки таблиц --- достаточно единожды загрузить диапазоны в программу, после чего интерфейс системы предоставляет мгновенные подсказки.
    
    \item \textbf{Гибкость настройки:} поддержка различных форматов диапазонов (включая популярные стандарты) позволяет легко переключаться между стратегиями для разных типов игр и покер-румов.
    
    \item \textbf{Обучающий эффект:} постоянное использование системы способствует неосознанному запоминанию диапазонов через регулярную практику, что особенно ценно для начинающих игроков.
\end{itemize}

Таким образом, разработанное ПО не только упрощает соблюдение стратегических принципов в процессе игры, но и служит инструментом для систематического обучения покерной стратегии, повышая дисциплинированность игрока и минимизируя влияние эмоциональных решений на результат игры.
